\documentclass[11pt]{article}
\usepackage[margin=1.1in]{geometry}
\usepackage{amsmath,amssymb,amsthm}
\usepackage{booktabs}
\usepackage{microtype}
\emergencystretch=2em
\usepackage[colorlinks=true,linkcolor=blue,citecolor=blue]{hyperref}

\newtheorem{theorem}{Theorem}[section]
\newtheorem{lemma}[theorem]{Lemma}
\newtheorem{proposition}[theorem]{Proposition}
\theoremstyle{remark}
\newtheorem{remark}[theorem]{Remark}

\newcommand{\dive}{\nabla\!\cdot\!}
\newcommand{\eps}{\varepsilon}
\newcommand{\B}{\mathcal{B}}

\title{The diffuse domain method on a mass-conservative adaptive
multigrid solver:\\ formulation, well-posedness, and
$\eps$-convergence analysis}
\author{DiffuseDomainBSAM documentation\\[2pt]
\small companion to \texttt{EllipticBSAM/docs/convergence\_theory.pdf}}
\date{}

\begin{document}
\maketitle

\begin{abstract}
We analyze the diffuse-domain reformulations solved by
\texttt{DiffuseDomainBSAM}: the flux/Neumann form of Feng, Guo,
Lowengrub \& Wise (2018), eq.~(6.18), and a Dirichlet penalty form, both
of which are instances of the general solver's equation
$-\dive(D_{dd}\nabla u)+C_{dd}u=f_{dd}$ with particular coefficients
built from a phase field $\phi$.  We prove uniform ellipticity of the
floored formulation and explain the formal matched-asymptotic prediction
of $O(\eps^2)$ for the smooth Neumann model under the required extension
and matching assumptions. This is not a uniform error theorem for every
coefficient, geometry, or penalty formulation. We also prove an
exact-reproduction lemma: for the analytic data extensions used in the
classical star-domain benchmark, the continuous diffuse-domain solution
\emph{satisfies the unfloored differential equation} for every $\eps$.
Equality of boundary-value solutions additionally requires compatible
outer data; the implemented floor is a separate perturbation. A
constant-normal extension gives measured rates near two in a small
circle study. We further quantify how the adaptive
hierarchy with
geometric band tagging resolves the layer at the finest resolution
while leaving the bulk coarse, and why the flux-balance corrections of
the underlying solver are essential at small $\eps$.
\end{abstract}

\section{Formulation}

Let $\Omega_1\subset\B$ be a smooth domain compactly contained in the
rectangle $\B$, described by a level set $\psi$ with $\psi<0$ in
$\Omega_1$, $\Gamma=\{\psi=0\}$.  Define the phase field
\begin{equation}\label{eq:phi}
  \phi(x) = \tfrac12\Bigl(1-\tanh\frac{3\psi(x)}{\eps}\Bigr),
\end{equation}
so $\phi\to 1$ in $\Omega_1$, $\phi\to0$ outside, with transition layer
of width $O(\eps)$.  Two reformulations of
\begin{equation}\label{eq:sharp}
  \beta u-\dive(D\nabla u)=f \text{ in }\Omega_1,
  \qquad\text{(BC on }\Gamma\text{)}
\end{equation}
onto $\B$ (homogeneous Neumann on $\partial\B$) are used:

\paragraph{Neumann data $D\,\partial_n u = g$ on $\Gamma$} (Feng et
al.\ eq.~(6.18); Li--Lowengrub--R\"atz--Voigt \cite{LLRV}):
\begin{equation}\label{eq:ddmn}
  \beta\,\phi\,u_\eps-\dive(\phi D\nabla u_\eps)
  = \phi f + |\nabla\phi|\,g
  \quad\text{in }\B,
\end{equation}
i.e.\ the general solver is called with
$D_{dd}=\phi D$, $C_{dd}=\beta\phi$, $f_{dd}=\phi f+|\nabla\phi|g$.

\paragraph{Dirichlet data $u=g$ on $\Gamma$} (penalty form):
\begin{equation}\label{eq:ddmd}
  -\dive(\phi D\nabla u_\eps)+\phi C u_\eps
  +\frac{\mu}{\eps^3}(1-\phi)(u_\eps-g)=\phi f ,
\end{equation}
i.e.\ $C_{dd}=\phi C+\mu\eps^{-3}(1-\phi)$,
$f_{dd}=\phi f+\mu\eps^{-3}(1-\phi)g$.

Here $f$ and $g$ denote smooth extensions of the data off $\Gamma$.

\section{Well-posedness and the $\phi$-floor}

In double precision $\tanh$ saturates and $\phi$ vanishes identically a
few layer-widths outside $\Gamma$, making \eqref{eq:ddmn} degenerate.
The implementation replaces $\phi$ by
$\phi_\delta=\max(\phi,\delta)$ in $D_{dd}$ (and in $C_{dd}$ for
\eqref{eq:ddmn}), $\delta=10^{-8}$ by default.

\begin{lemma}[continuous coercivity at fixed floor]
Suppose $0<D_{\min}\le D\le D_{\max}$ and $0<\delta<1$.
Then $D_{dd}\ge\delta D_{\min}>0$. For the Neumann formulation with
$\beta>0$, $C_{dd}\ge\beta\delta>0$, and its weak bilinear form is
coercive on $H^1(\B)$, hence has a unique solution for bounded data.
For the penalty formulation assume $C\ge0$, $\mu>0$, and a positive-measure
exterior region; its positive reaction controls the constant mode and a
weighted Poincare inequality gives the same fixed-parameter conclusion.
If $\beta=0$ with only Neumann outer data, compatibility and a mean
normalization are required. These continuous statements do not establish
that an assembled QIF composite matrix is an $M$-matrix: interpolation,
coefficient ratios, and mesh structure must also be checked.
\end{lemma}

\begin{lemma}[effect of the floor]\label{lem:floor}
Let $0<\delta<1/2$. Let $u_\eps$ solve \eqref{eq:ddmn}, and let
$u_\eps^\delta$ solve the equation with diffusion $\phi_\delta D$ and
reaction $\beta\phi_\delta$, keeping the assembled right-hand side
$f_{dd}=\phi f+|\nabla\phi|g$ unchanged, as in the implementation.
The coefficient perturbation is supported in
$\{\phi<\delta\}=\{\psi>\psi_\delta\}$ with
$\psi_\delta=\tfrac{\eps}{6}\ln\frac{1-\delta}{\delta}
\approx\tfrac{\eps}{6}\ln\frac1\delta$, a logarithmic distance in
layer units outside $\Gamma$. Assume $\beta>0$, homogeneous outer Neumann data,
and the stated ellipticity bounds on $D$. Testing the difference equation with
the difference and using $\phi\le\phi_\delta\le\phi+\delta$,
\[
  \|\sqrt{\phi_\delta}\,\nabla(u_\eps^\delta-u_\eps)\|_{L^2(\B)}^2
  +\beta\|\sqrt{\phi_\delta}\,(u_\eps^\delta-u_\eps)\|_{L^2(\B)}^2
  \;\le\; C\,\delta\,
  \bigl(\|\nabla u_\eps\|^2_{L^2(\phi<\delta)}
  +\|u_\eps\|^2_{L^2(\phi<\delta)}\bigr),
\]
Since $\phi\ge1/2$ inside $\Omega_1$, this controls the interior
$H^1$ error by $O(\delta^{1/2})$ times the displayed exterior norms.
Uniformity in $\eps$ requires an independent bound on those norms.
The estimate is an energy estimate, not an $L^\infty$ bound, and it
does not certify that a particular floor is below a requested tolerance.
\end{lemma}

\section{Sharp-interface limit: matched asymptotics}\label{sec:asymp}

We expand \eqref{eq:ddmn}.  Assume first that $\psi=d$, the
\emph{signed distance} to $\Gamma$ ($|\nabla d|=1$), let $\kappa$ be
the curvature, and use normal--tangential coordinates $(d,s)$ with the
stretched variable $z=d/\eps$ and profile
$\Phi(z)=\tfrac12(1-\tanh 3z)$, so that
\[
  \phi=\Phi(z),\qquad
  |\nabla\phi|=\eps^{-1}|\Phi'(z)|,\qquad
  \int_{-\infty}^{\infty}|\Phi'(z)|\,dz=1 .
\]
Write outer expansions $u_\eps=u^0+\eps u^1+\cdots$ in $\Omega_1$ (and
in $\B\setminus\overline\Omega_1$) and an inner expansion
$\tilde u(z,s)=\tilde u^0+\eps\tilde u^1+\eps^2\tilde u^2+\cdots$.  In
the layer,
\[
  \dive(\phi D\nabla u)=
  \eps^{-2}\partial_z\bigl(\Phi D\,\partial_z\tilde u\bigr)
  +\eps^{-1}\kappa\,\Phi D\,\partial_z\tilde u
  +O(1).
\]

\paragraph{Order $\eps^{-2}$.}
$\partial_z(\Phi D\,\partial_z\tilde u^0)=0$, so
$\Phi D\,\partial_z\tilde u^0=a(s)$.  Matching to the interior outer
solution as $z\to-\infty$ ($\Phi\to1$, $\partial_z\tilde u^0\to0$)
forces $a\equiv0$; since $\Phi>0$ for all finite $z$,
$\partial_z\tilde u^0\equiv 0$:
\begin{equation}\label{eq:u0const}
  \tilde u^0=\tilde u^0(s)=u^0\big|_\Gamma(s),
\end{equation}
the leading inner solution is constant across the layer.  The interior
outer problem at $O(1)$ is the sharp equation \eqref{eq:sharp}.

\paragraph{Order $\eps^{-1}$ --- recovery of the Neumann condition.}
Using \eqref{eq:u0const},
\[
  -\partial_z\bigl(\Phi D\,\partial_z\tilde u^1\bigr)
  = |\Phi'(z)|\,g(s) .
\]
Integrating over $z\in\mathbb R$ and using
$\Phi D\partial_z\tilde u^1\to D\,\partial_d u^0|_\Gamma$ as
$z\to-\infty$ (matching) and $\to0$ as $z\to+\infty$ (boundedness;
verified below), and $\int|\Phi'|=1$:
\begin{equation}\label{eq:bc}
  D\,\partial_n u^0 = g \quad\text{on }\Gamma .
\end{equation}
Thus $u^0$ solves the sharp problem.  Solving for the profile,
\[
  \Phi D\,\partial_z\tilde u^1(z)
  = D\partial_n u^0\big|_\Gamma - g\!\int_{-\infty}^{z}\!|\Phi'|
  = g\!\int_{z}^{\infty}\!|\Phi'|\,dz' ,
\]
whose right side decays like $e^{-6z}$, matching $\Phi\sim e^{-6z}$, so
$\partial_z\tilde u^1$ is bounded: the expansion is consistent.

\paragraph{Next order and scope of the formal prediction.}
For the smooth, constant-coefficient Neumann model with signed distance
and constant-normal extensions, the full matching calculation in
\cite{LervagLowengrub} gives $u^1=0$ when $\beta>0$; in the pure
Neumann case the constant correction must be fixed by a consistent gauge.
The symmetric profile satisfies
$\int_{\mathbb R}z|\Phi'(z)|\,dz=0$ and
$\int_{\mathbb R}(\Phi(z)-\mathbf 1_{z<0})\,dz=0$.
These identities help cancel the first correction; by themselves they
are not a substitute for all solvability and matching conditions.
The resulting formal prediction is
\begin{equation}\label{eq:eps2}
  \|u_\eps-u\|_{L^2(\Omega_1)} = O(\eps^2),
\end{equation}
under those assumptions. Extending it to variable $D$, an arbitrary data
extension, a finite outer boundary, the floor, and a changing adaptive
mesh requires additional error estimates.

\begin{lemma}[exact reproduction by compatible extensions]\label{lem:exact}
Let $u_e\in C^2(\overline\B)$ and suppose the data extensions satisfy,
pointwise in $\B$,
\[
  \beta u_e-\dive(D\nabla u_e)=f
  \qquad\text{and}\qquad
  |\nabla\phi|\,g=-\nabla\phi\cdot D\nabla u_e .
\]
Then $u_e$ satisfies the unfloored differential equation \eqref{eq:ddmn}
\emph{exactly, for every} $\eps$: indeed
$\beta\phi u_e-\dive(\phi D\nabla u_e)
=\phi(\beta u_e-\dive(D\nabla u_e))-\nabla\phi\cdot D\nabla u_e
=\phi f+|\nabla\phi|g$.
The second condition holds whenever $g$ is extended as
$g=D\nabla u_e\cdot\nabla\psi/|\nabla\psi|$, because
$\nabla\phi=\phi'(3\psi/\eps)\tfrac{3}{\eps}\nabla\psi$ with $\phi'<0$,
so $|\nabla\phi|=-\phi'\tfrac3\eps|\nabla\psi|$.
In particular, for the test data $u_e=(x^2+y^2)/4$,
$f=u_e-\Delta u_e=r^2/4-1$ and $g=\nabla u_e\cdot n_\psi$ (any
geometry, including the flower), this is an exact differential identity.
To be the solution on $\B$, $u_e$ must also satisfy the imposed outer
boundary condition. The quadratic does not satisfy homogeneous Neumann
conditions on a finite box, and the floor changes the identity. Thus
the benchmark is dominated by layer discretization when these other
errors are small, but does not have identically zero total model error.
This observation also applies to the corresponding test construction of
Feng et al.\ (their Sec.~6.4): the $\eps$-convergence reported there is
convergence of the \emph{numerical} solution under the coupled
$(\eps,h)$ schedule.  To probe the genuine $O(\eps^2)$ \emph{model}
convergence one must use an extension that breaks the identity --- the
canonical choice being the constant-in-normal extension
$g\equiv g|_\Gamma$ --- and separate the discretization error, e.g.\ by
Richardson extrapolation in $h$ at fixed $\eps$ (test D1b).
\end{lemma}

\begin{remark}[non-distance level sets]\label{rem:nondist}
If $\Gamma$ is described by a level set $\psi$ with
$|\nabla\psi|\ne1$ (e.g.\ the flower $\psi=r-R_0(\theta)$ of Feng et
al.\ eq.~(6.19)), then in the true normal coordinate
$\psi = |\nabla\psi|_\Gamma\, d + \tfrac12 (\partial^2_n\psi)\, d^2 +
O(d^3)$, and $\phi=\Phi(\psi/\eps)$ has a spatially varying physical
width and can be skewed in $d$. The signed-distance calculation above
therefore cannot be transferred without re-expanding the geometry and
data. No universal rate between one and two follows from
$|\nabla\psi|\ne1$ alone: even a constant rescaling of a distance function
changes the width without destroying symmetry. Reinitialization can
simplify the analysis, but does not itself prove a convergence rate.
\end{remark}

\begin{remark}[corners]\label{rem:corners}
For a polygonal $\Omega_1$ with the exact signed distance as level set,
$|\nabla\psi|=1$ a.e.\ but $\nabla\psi$ jumps across the corner
bisectors.  Two distinct effects follow.  (i)~In the \emph{flux} form
\eqref{eq:ddmn} the extended datum $g=\nabla u\cdot n_\psi$ inherits the
jump: $f_{dd}$ is discontinuous along five lines (for a pentagon)
inside the band, and the corner neighborhoods converge at reduced order
in $L^2$ (measured: rate $\approx0.7$ in $L^2$, $\approx1.7$ in
$L^\infty$, at fixed $\eps/h$) in the recorded test. Discontinuous extended data can reduce regularity;
these observations alone do not isolate the source of the measured rate.  (ii)~The \emph{penalty} form
\eqref{eq:ddmd} never references the normal; corners enter only through
the geometry of $\phi$, and the measured $\eps$-rates on the pentagon
are $1.5$--$1.8$ over the recorded range. The smooth Neumann
expansion does not establish the asymptotic penalty rate at corners.
\end{remark}

\paragraph{Dirichlet penalty.}  For \eqref{eq:ddmd} the dominant inner
balance at $O(\eps^{-3})$ is $\mu(1-\Phi)(\tilde u^0-g)=0$; since
$1-\Phi>0$ for all finite $z$, $\tilde u^0\equiv g$ throughout the
layer and matching gives $u^0=g$ on $\Gamma$: the Dirichlet condition
is enforced at leading order.  The $\eps^{-3}$ scaling makes the
penalty dominate the diffusion ($\eps^{-2}$) by one order, and the
leading-order argument does not determine the error rate. The paper
\cite{LervagLowengrub} analyzes Neumann and Robin approximations and is
not a proof of second-order accuracy for this Dirichlet penalty.
The disc suite measures rates $1.8$, $2.1$, and $3.6$ over its finite
schedule; these are observations for that problem, not a general
second-order theorem.

\section{Interaction with the adaptive solver}

\begin{remark}[band tagging and actual coverage]\label{lem:band}
Tag every cell with $|\psi|<w\eps$ (TagWidth $w=2.5$ by default) on
every level, with buffer $b\ge2$ cells.  Then outside the tagged band
$|\phi-\mathbf 1_{\psi<0}|\le\tfrac12(1-\tanh 3w)\approx 3\cdot10^{-7}$
($w=2.5$). This is a phase-field tail bound, not a guarantee of
finest-level coverage: tags are sampled at cell centers, clustering may
drop thin fragments, and the level cap may stop refinement. Check the
visible cells intersecting the band. For narrow layers, a useful
level-dependent tag width is $\max(w\eps,3h_k)$ for signed distance;
this reduces missed coarse-grid fragments, but still needs coverage
verification. The tail is nonzero even outside the tagged band.
\end{remark}

\begin{remark}[a local coefficient-ratio check]
For two sample points separated by at most $h$, assume
$|\nabla\psi|\le G$ along the connecting segment. Since
$|\partial_\psi\log\phi|\le6/\eps$, the phase-field ratio is bounded
below by $e^{-6Gh/\eps}$; flooring preserves this bound. If the physical
diffusivity ratio is bounded below by $\rho_D>0$, a sufficient check
for a local ratio $D_2/D_1\ge1/5$ is
\[
  \rho_D e^{-6Gh/\eps}\ge\tfrac15.
\]
For constant physical diffusivity and signed distance this reduces to
$\eps/h\ge6/\log5\approx3.73$. Use the actual spacing and geometry at
the coarse--fine stencil, not just the finest spacing. This local check
does not certify diagonal dominance or positivity of the entire
composite QIF matrix.
\end{remark}

\begin{remark}[why the flux-balance corrections matter here]
The DDM coefficients are exponentially small outside the band, so the
solution is determined almost entirely through the layer --- the region
crossed by every coarse--fine interface.  Any flux imbalance there acts
like a spurious source of relative size $O(h^2)$ \emph{concentrated
exactly where all the physics happens}.  The suite's ablation shows the
consequence at $\eps=0.1$ (4-level mesh): disabling the corrections
inflates the interior $L^2$ error from $3.0\cdot10^{-6}$ to
$5.2\cdot10^{-3}$ (a factor $\sim$1700) and the mass defect from
$10^{-13}$ to $2\cdot10^{-3}$ --- the multilevel analogue of the
paper's Table~6, and the mechanism behind the divergence of earlier
in-house attempts at $\ge3$ levels.
\end{remark}

\section{Numerical verification}\label{sec:num}

Produced by \texttt{tests/runDdTests.m}; complete logs and tables in
\texttt{TEST\_REPORT.md}.  Highlights (interior $L^2$ errors against the
sharp solution):

\begin{itemize}
\item \textbf{Flower (paper Sec.~6.4 replication), Neumann form,
analytic extensions:} Lemma~\ref{lem:exact} explains the differential
cancellation, subject to its floor and outer-boundary qualifications.
The $(\eps,h)$-schedule errors are dominated by layer discretization,
decaying with observed rates $\approx1.5$ in $\eps$ at fixed
$\eps/h=20.5$ down to $1.1\cdot10^{-6}$ at $\eps=0.05$ (effective
$2048^2$); the adaptive 4-level solution agrees with the uniform
solution of equal finest resolution to $\lesssim1\%$ while using
$5$--$10\times$ fewer cells.
\item \textbf{Circle, constant-normal extension (model-error diagnostic):}
The diagnostic $(4E(h/2)-E(h))/3$ assumes a smooth $h^2$ expansion of
the error norm; it estimates a model-error floor but is not a certified
extrapolation of the solution. Its measured $\eps$-rate is near two,
consistent with \eqref{eq:eps2}. The displayed ratios of total errors
at fixed $\eps$ are near zero order because model error dominates;
they do not measure spatial order.
\item \textbf{Disc, Dirichlet penalty:} rates $1.8$--$2.1$ (with a
superconvergent last interval as the discretization error crosses the
model error).
\item \textbf{Conservation:} relative mass defect
$10^{-16}$--$10^{-13}$ at all $\eps$ with the corrections on; up to
$10^{-3}$ with them off, accompanied by a $\sim10^3\times$ loss of
interior accuracy at $\eps\le0.1$ on the 4-level meshes.
\item \textbf{Solver:} V(3,3) factors $0.05$--$0.16$ per cycle on the
adaptive DDM hierarchies, independent of $\eps$ despite coefficient
degeneracy of eight orders of magnitude.
\end{itemize}

\begin{thebibliography}{9}
\bibitem{FGLW2018} W.~Feng, Z.~Guo, J.~S.~Lowengrub, S.~M.~Wise,
A mass-conservative adaptive FAS multigrid solver for cell-centered
finite difference methods on block-structured, locally-cartesian grids,
\emph{J.~Comput.~Phys.} 352 (2018) 463--497.
\bibitem{LLRV} X.~Li, J.~Lowengrub, A.~R\"atz, A.~Voigt, Solving PDEs in
complex geometries: a diffuse domain approach, \emph{Commun.\ Math.\
Sci.} 7 (2009) 81--107.
\bibitem{LervagLowengrub} K.~Y.~Lerv\r{a}g, J.~Lowengrub, Analysis of
the diffuse-domain method for solving PDEs in complex geometries,
\emph{Commun.\ Math.\ Sci.} 13 (2015) 1473--1500.
\end{thebibliography}

\end{document}
